\chapter{Angular}

Abbiamo visto che una Single Page Application comporta una certa complessità:
routing lato client, generazione delle viste, componenti e riuso del codice. Un
\term{framework di front-end} è un insieme predefinito di componenti software,
strumenti e buone pratiche che fa da fondamento per lo sviluppo di SPA.

\section{Che cos'è Angular}

\term{Angular} è uno dei framework di front-end più diffusi, in particolare in
Italia e per le applicazioni aziendali. È sviluppato da Google; la versione 2.0
è stata rilasciata nel 2016 e in questo corso si usa la \textbf{versione 17}.

Le sue caratteristiche distintive:

\begin{itemize}
  \item forte attenzione agli \textbf{strumenti} per lo sviluppatore e alla
        produttività;
  \item \textbf{tutto incluso}: Angular è un framework completo e
        \textit{opinionated} (include il meccanismo di iniezione delle
        dipendenze, il routing e altro), il che riduce l'affaticamento
        decisionale e garantisce coerenza fra progetti;
  \item \textbf{supporto alla migrazione}: strumenti dedicati assistono nel
        passaggio a versioni più recenti;
  \item \textbf{buone pratiche fin dall'inizio}: impone l'uso di TypeScript.
\end{itemize}

\section{Primi passi}

\begin{lstlisting}[style=shell]
npm install -g @angular/cli
ng version
\end{lstlisting}

Come ambiente si usa Visual Studio Code; è consigliata (ma facoltativa)
l'estensione \textit{Angular Language Service}.

\begin{lstlisting}[style=shell]
ng new angular-app --create-application
\end{lstlisting}

Il comando chiede quale formato di foglio di stile usare (SCSS) e se abilitare
il rendering lato server, poi genera lo scheletro del progetto. Per avviarlo:

\begin{lstlisting}[style=shell]
cd angular-app
ng serve
\end{lstlisting}

\begin{console}[Terminale]
Initial Chunk Files | Names     |  Raw Size
polyfills.js        | polyfills |  82.71 kB
main.js             | main      |  23.23 kB
styles.css          | styles    |     96 bytes
                    | Initial Total | 106.03 kB

Application bundle generation complete. [1.176 seconds]
Watch mode enabled. Watching for file changes...
  ->  Local:   http://localhost:4200/
\end{console}

\section{I componenti}

I \term{componenti} sono i mattoni delle applicazioni Angular. Ogni componente
è responsabile di definire funzione e aspetto di un elemento:

\begin{itemize}
  \item il proprio codice (stato, logica di business, eventi);
  \item il proprio layout HTML;
  \item i propri stili CSS.
\end{itemize}

In Angular un componente può contenere altri componenti: si può pensare a
un'applicazione Angular come a un \textbf{albero di componenti}.

\begin{figure}[H]
  \centering
  \begin{tikzpicture}[font=\Lato\scriptsize,
      level distance=1.1cm,
      n/.style={rounded corners=3pt, draw=primary!60!black, line width=0.7pt,
                fill=primary!10!white, inner sep=4pt, align=center},
      edge from parent/.style={draw=neutral!50!black, line width=0.6pt,
                               edge from parent path=
                               {(\tikzparentnode.south) -- (\tikzchildnode.north)}},
      level 1/.style={sibling distance=4.4cm},
      level 2/.style={sibling distance=2.6cm}]
    \node[n] {Pagina lista}
      child { node[n] {Form nuovo\\ elemento} }
      child { node[n] {Lista}
        child { node[n] {Elemento} }
        child { node[n] {Elemento} } };
  \end{tikzpicture}
  \caption{Un'applicazione Angular è un albero di componenti.}
  \label{fig:angular-tree}
\end{figure}

\subsection{Anatomia di un componente}

Un componente è sostanzialmente una \textbf{classe TypeScript}, annotata con
metadati passati come oggetto al \term{decoratore} \code{@Component}. I metadati
specificano, fra le altre cose, il layout HTML e gli stili.

\begin{affianco}[0.5][c]
\codicepiccolo
\begin{lstlisting}[style=ts, name=app.component.ts]
import { Component } from '@angular/core';

@Component({
  selector: 'app-root',
  standalone: true,
  template: "<h1>Hello Angular!</h1>",
  styles: "h1 { font-family: sans-serif; color: darkblue; }",
})
export class AppComponent {}
\end{lstlisting}
\risultato
\begin{browser}[AngularApp]
  {\Lato\bfseries\LARGE\color{css-darkblue} Hello Angular!\par}
\end{browser}
\end{affianco}

Il \code{selector} dice ad Angular \textbf{dove} disegnare il componente:
Angular ne crea un'istanza per ogni elemento HTML corrispondente.

\begin{affianco}[0.42][c]
\begin{lstlisting}[style=html, name=index.html]
<!doctype html>
<html lang="it">
  <head>
    <meta charset="utf-8">
    <title>AngularApp</title>
  </head>
  <body>
    <app-root></app-root>
  </body>
</html>
\end{lstlisting}
\risultato
\begin{immagine}[Angular trova \code{<app-root>} e vi disegna\\
                 il template del componente.]
\begin{tikzpicture}
  \node[mdom] (b) at (0,0) {<body>};
  \node[mnuovo] (a) at (0,-0.9) {<app-root>};
  \node[mdom] (h) at (0,-1.8) {<h1>};
  \node[mtesto] (t) at (0,-2.7) {Hello Angular!};
  \draw[mramo] (b.south) -- (a.north); \draw[mramo] (a.south) -- (h.north);
  \draw[mramo] (h.south) -- (t.north);
  \node[mnota, anchor=west] at ([xshift=10pt]a.east) {istanza di \code{AppComponent}};
  \draw[wtbrace] ([xshift=10pt]t.south east) -- ([xshift=10pt]t.south east |- h.north)
    node[mnota, midway, right=5pt] {dal \code{template}};
\end{tikzpicture}
\end{immagine}
\end{affianco}

Template e stili si possono anche mettere in file separati, con percorsi
relativi alla posizione della classe TypeScript:

\begin{affianco}[0.46][c]
\begin{lstlisting}[style=ts, name=app.component.ts]
@Component({
  selector: 'app-root',
  standalone: true,
  templateUrl: './app.component.html',
  styleUrl: './app.component.scss'
})
export class AppComponent {}
\end{lstlisting}
\risultato
\begin{immagine}
\begin{tikzpicture}[
    f/.style={wtnodeplain, font=\FiraCodeNL\scriptsize}]
  \node[f] (ts) at (0,0) {app.component.ts};
  \node[f, anchor=west] (ht) at (3.3,0.55) {app.component.html};
  \node[f, anchor=west] (sc) at (3.3,-0.55) {app.component.scss};
  \draw[mrif] (ts.east) -- node[mcod, above, sloped]{templateUrl} (ht.west);
  \draw[mrif] (ts.east) -- node[mcod, below, sloped]{styleUrl} (sc.west);
  \node[mnota, anchor=south, font=\FiraCodeNL\scriptsize\bfseries]
    at ([yshift=8pt]ht.north -| 2.3,0) {src/app/};
\end{tikzpicture}
\end{immagine}
\end{affianco}

\subsection{Lo stato e l'interpolazione}

\begin{affianco}[0.5][c]
\codicepiccolo
\begin{lstlisting}[style=ts, name=app.component.ts]
@Component({
  selector: 'app-root',
  standalone: true,
  template: "<h1>Hello {{name.toUpperCase()}}!</h1>",
  styles: "h1 { font-family: sans-serif; color: darkblue; }",
})
export class AppComponent {
  name = 'Web Technologies';
}
\end{lstlisting}
\risultato
\begin{browser}[AngularApp]
  {\Lato\bfseries\LARGE\color{css-darkblue} Hello WEB TECHNOLOGIES!\par}
\end{browser}
\wtnota{\code{name.toUpperCase()} è valutata e il risultato\\inserito nel titolo.}
\end{affianco}

Le espressioni fra \code{\{\{ \}\}} nei template Angular vengono valutate, e il
risultato (convertito in stringa) viene inserito nel template. È lo stesso
principio dei motori di template visti nel capitolo 11, ma applicato
\textbf{lato client} e ricalcolato automaticamente a ogni cambiamento di stato.

\subsection{Comporre componenti}

\begin{lstlisting}[style=shell]
ng generate component counter
\end{lstlisting}

\begin{console}[Terminale]
CREATE src/app/counter/counter.component.html (22 bytes)
CREATE src/app/counter/counter.component.spec.ts (603 bytes)
CREATE src/app/counter/counter.component.ts (239 bytes)
CREATE src/app/counter/counter.component.scss (0 bytes)
\end{console}

\begin{affianco}[0.52][c]
\codicepiccolo
\begin{lstlisting}[style=ts, name=app.component.ts]
import { Component } from '@angular/core';
import { CounterComponent } from './counter/counter.component';

@Component({
  selector: 'app-root',
  standalone: true,
  imports: [CounterComponent],
  template: `
    <h1>Hello {{name.toUpperCase()}}!</h1>
    <app-counter />
  `,
  styles: "h1 { font-family: sans-serif; color: darkblue; }",
})
export class AppComponent {
  name = 'Web Technologies';
}
\end{lstlisting}
\risultato
\begin{browser}[AngularApp]
  {\Lato\bfseries\LARGE\color{css-darkblue} Hello WEB TECHNOLOGIES!\par}\vspace{6pt}
  \tikz\node[draw=accent!60!black, dashed, rounded corners=2pt, inner sep=5pt,
    label={[font=\FiraCodeNL\tiny, text=accent!70!black, inner sep=1pt,
            anchor=south east]north east:<app-counter>}]
    {counter works!};
\end{browser}
\wtnota{Il template generato da \code{ng generate} contiene soltanto\\
        \code{<p>counter works!</p>}.}
\end{affianco}

\info{Componenti standalone}{
  I componenti \textit{standalone} hanno \code{standalone: true} nei metadati e
  possono importare direttamente altri componenti, direttive e \textit{pipe}
  usati nei propri template. Il codice Angular precedente alla versione 14
  usava un meccanismo diverso, chiamato \code{NgModule}. In questo corso si
  usano soltanto componenti standalone.}

\section{Gestione degli eventi}

In Angular gli eventi si legano ai gestori con la sintassi a \textbf{parentesi
tonde}:

\begin{lstlisting}[style=html]
<button (click)="increment()">Click me</button>
\end{lstlisting}

dove \code{increment()} è un metodo del componente.

\begin{affianco}[0.54][c]
\codicepiccolo
\begin{lstlisting}[style=ts, name=counter.component.ts]
@Component({
  selector: 'app-counter',
  standalone: true,
  template: `
    <button (click)="increment()">Click me</button>
    <p>The button was clicked {{counter}} time(s).</p>
  `,
  styles: 'p, button { font-size: 1.25em; font-family: sans-serif; }'
})
export class CounterComponent {
  counter: number = 0;

  increment() {
    this.counter = this.counter + 1;
  }
}
\end{lstlisting}
\risultato
\begin{browser}[AngularApp]
  {\Lato\bfseries\Large\color{css-darkblue} Hello WEB TECHNOLOGIES!\par}\vspace{6pt}
  \wtbottone{{\normalsize\Lato Click me}}\par\vspace{6pt}
  {\Lato\large The button was clicked 3 time(s).\par}
\end{browser}
\begin{immagine}[Nessuna riga di codice tocca il DOM.]
\begin{tikzpicture}[
    s/.style={wtnodeplain, font=\FiraCodeNL\scriptsize},
    a/.style={-{Stealth[length=4pt,width=3.5pt]}, line width=0.75pt,
      draw=neutral!45!black}]
  \node[s, font=\Lato\scriptsize] (c) at (0,0) {clic};
  \node[s] (i) at (1.75,0) {increment()};
  \node[s] (v) at (4.3,0) {this.counter = 3};
  \node[wtnodeok, font=\Lato\scriptsize, inner sep=3pt] (dom) at (4.3,-1.05)
    {Angular aggiorna \code{\{\{counter\}\}}};
  \draw[a] (c) -- (i); \draw[a] (i) -- (v);
  \draw[a] (v) -- node[mnota, right]{da sé} (dom);
\end{tikzpicture}
\end{immagine}
\end{affianco}

Si noti che non c'è alcun codice che aggiorna il DOM: modificando
\code{this.counter}, Angular ridisegna da sé la parte di template che dipende da
quel valore. È esattamente il lavoro manuale che nel capitolo 18 dovevamo
scrivere a mano.

\section{Controllo di flusso nei template}

Spesso serve disegnare condizionalmente una parte di template, o iterare su più
elementi. Angular supporta questi scenari con la sintassi \code{@if}/\code{@else}
e \code{@for}. Prima di Angular 17 si usavano le direttive \code{ngIf} e
\code{ngFor}.

\begin{affianco}[0.46][c]
\begin{lstlisting}[style=html, name=counter.component.html]
<button (click)="increment()">Click me</button>
<p>
  The button was clicked {{counter}}
  @if (counter === 1) { time } @else { times }.
</p>
\end{lstlisting}
\risultato
\begin{browser}[counter = 1]
  \wtbottone{{\normalsize\Lato Click me}}\par\vspace{4pt}
  {\Lato\large The button was clicked 1 time.\par}
\end{browser}
\begin{browser}[counter = 2]
  \wtbottone{{\normalsize\Lato Click me}}\par\vspace{4pt}
  {\Lato\large The button was clicked 2 times.\par}
\end{browser}
\end{affianco}

\begin{esempio}{Iterare con \code{@for}}
Aggiungiamo al contatore la registrazione dei momenti in cui il pulsante è
stato premuto.

\begin{lstlisting}[style=ts, name=counter.component.ts]
export class CounterComponent {
  counter: number = 0;
  clickTimestamps: Date[] = [];

  increment() {
    this.counter = this.counter + 1;
    this.clickTimestamps.push(new Date());
  }
}
\end{lstlisting}

\begin{affianco}[0.4][c]
\codicepiccolo
\begin{lstlisting}[style=html, name=counter.component.html]
<ul>
  @for (ts of clickTimestamps; track ts) {
    <li>{{ts}}</li>
  } @empty {
    <li>Il pulsante non e' mai stato premuto!</li>
  }
</ul>
\end{lstlisting}
\risultato
\begin{browser}[nessun clic]
  \begin{itemize}[leftmargin=14pt, itemsep=0pt, topsep=0pt]
    \item Il pulsante non e' mai stato premuto!
  \end{itemize}
\end{browser}
\begin{browser}[dopo due clic]
  \footnotesize
  \begin{itemize}[leftmargin=14pt, itemsep=0pt, topsep=0pt]
    \item Mon Oct 05 2026 10:15:32 GMT+0200 (Ora legale dell'Europa centrale)
    \item Mon Oct 05 2026 10:15:34 GMT+0200 (Ora legale dell'Europa centrale)
  \end{itemize}
\end{browser}
\wtnota{Interpolata, una \code{Date} diventa la stringa di \code{toString()}.}
\end{affianco}

La clausola \code{track} indica ad Angular come identificare univocamente gli
elementi della collezione: serve a ridisegnare solo le voci effettivamente
cambiate invece dell'intera lista. È la soluzione al problema del
\guillemotleft rendering efficiente\guillemotright\ rimasto aperto nel capitolo 18.
\end{esempio}

\section{Comunicazione fra componenti}

A volte i componenti devono scambiarsi dati: un componente padre passa
informazioni ai figli (per configurarne il comportamento o per passare dati da
visualizzare), oppure un figlio deve comunicare al padre che è accaduto
qualcosa. In Angular questi schemi si realizzano con i decoratori
\code{@Input()} e \code{@Output()}.

\subsection{Dal padre al figlio: @Input}

Nel componente figlio, le proprietà che il padre può impostare si decorano con
\code{@Input()}. Può esserci un valore predefinito, usato quando il padre non ne
fornisce uno.

\begin{affianco}[0.6][c]
\begin{lstlisting}[style=ts, name=counter.component.ts]
import { Component, Input } from '@angular/core';

export class CounterComponent {
  @Input() buttonText = "Click me";
  counter: number = 0;
  clickTimestamps: Date[] = [];

  increment() {
    this.counter = this.counter + 1;
    this.clickTimestamps.push(new Date());
  }
}
\end{lstlisting}
\risultato
\begin{lstlisting}[style=html, name=counter.component.html]
<button (click)="increment()">
  {{buttonText}}
</button>
\end{lstlisting}
\end{affianco}

Il padre passa il dato impostando una proprietà sull'elemento:

\begin{affianco}[0.46][c]
\begin{lstlisting}[style=ts, name=app.component.ts]
@Component({
  selector: 'app-root',
  standalone: true,
  imports: [CounterComponent],
  template: `
    <h1>Hello {{name.toUpperCase()}}!</h1>
    <app-counter buttonText="CLICK HERE"/>
  `,
})
export class AppComponent {
  name = 'Web Technologies';
}
\end{lstlisting}
\risultato
\begin{browser}[AngularApp]
  {\Lato\bfseries\LARGE\color{css-darkblue} Hello WEB TECHNOLOGIES!\par}\vspace{6pt}
  \wtbottone{{\normalsize\Lato CLICK HERE}}
\end{browser}
\end{affianco}

\begin{figure}[H]
  \centering
  \begin{tikzpicture}[font=\Lato\scriptsize,
      lab/.style={font=\Lato\scriptsize, align=center, inner sep=2pt}]
    \node[font=\FiraCodeNL\normalsize] (c) at (0,0)
      {<app-counter buttonText="CLICK HERE"/>};

    \node[lab, text=accent!70!black, anchor=south east] (l1) at (-1.3,0.9)
      {selettore del figlio,\\ nel template del padre};
    \node[lab, text=codeattr, anchor=north] (l2) at (0.35,-0.65)
      {proprietà \code{@Input()}\\ del figlio};
    \node[lab, text=codestring, anchor=north west] (l3) at (2.0,-0.65)
      {valore passato\\ dal padre};

    \draw[wtdash] (l1.south east) -- (-1.5,0.25);
    \draw[wtdash] (l2.north) -- (0.35,-0.25);
    \draw[wtdash] (l3.north west) -- (2.2,-0.25);
  \end{tikzpicture}
  \caption{Anatomia del passaggio di dati dal padre al figlio.}
  \label{fig:angular-input}
\end{figure}

Questo schema è molto diffuso nei framework di front-end: in React e Vue lo
stesso meccanismo si chiama \term{props}.

\subsection{Property binding}

Angular permette di impostare \textbf{dinamicamente} i valori degli attributi di
elementi HTML e componenti, legandoli a proprietà del componente corrente. Il
nome dell'attributo da legare si scrive fra parentesi quadre, e il valore è il
nome della proprietà.

\begin{affianco}[0.46][c]
\begin{lstlisting}[style=html, name=app.component.html]
<h1>Hello {{name.toUpperCase()}}!</h1>
<app-counter [buttonText]="btnMsg"/>
<img [src]="imageUrl"/>
\end{lstlisting}

\begin{lstlisting}[style=ts, name=app.component.ts]
export class AppComponent {
  name = 'Web Technologies';
  btnMsg = "CLICK HERE";
  imageUrl = "https://picsum.photos/300/150";
}
\end{lstlisting}
\risultato
\begin{browser}[AngularApp]
  {\Lato\bfseries\LARGE\color{css-darkblue} Hello WEB TECHNOLOGIES!\par}\vspace{6pt}
  \wtbottone{{\normalsize\Lato CLICK HERE}}\par\vspace{6pt}
  \tikz{\fill[neutral!12!white] (0,0) rectangle (4,2);
    \fill[green-gradient-6!45!white] (0,0) -- (1.2,1.1) -- (2,0.55) -- (2.9,1.35)
      -- (4,0.3) -- (4,0) -- cycle;
    \fill[yellow-gradient-5] (3.2,1.55) circle (0.22);}
\end{browser}
\wtnota{L'immagine casuale $300\times150$ indicata da \code{imageUrl}.}
\end{affianco}

\warning{Property binding contro interpolazione}{
  Ci si potrebbe chiedere che differenza ci sia fra \code{<img [src]="imageUrl"/>}
  e \code{<img src="\{\{imageUrl\}\}"/>}. In quell'esempio il risultato è lo
  stesso, ma non è sempre così: l'\textbf{interpolazione} valuta l'espressione e
  ne \textbf{converte il risultato in stringa}. Per passare un \textbf{oggetto}
  o un \textbf{array} come proprietà di un elemento bisogna necessariamente
  usare il property binding.}

\subsection{Dal figlio al padre: @Output}

Per creare un canale di comunicazione dal figlio al padre si usa il decoratore
\code{@Output} su una proprietà di classe, assegnandole un valore di tipo
\code{EventEmitter}. Per inviare un evento sul canale, il figlio invoca il
metodo \code{emit()}.

\begin{lstlisting}[style=ts, name=counter.component.ts]
export class CounterComponent {
  @Input() buttonText = "Click me";
  @Output() counterClickedTenTimes = new EventEmitter<number>();

  counter: number = 0;
  clickTimestamps: Date[] = [];

  increment() {
    this.counter++;
    this.clickTimestamps.push(new Date());

    if (this.counter % 10 === 0)
      this.counterClickedTenTimes.emit(this.counter);
  }
}
\end{lstlisting}

Il padre reagisce definendo un gestore, come per un evento qualunque:

\begin{affianco}[0.48][c]
\codicepiccolo
\begin{lstlisting}[style=html, name=app.component.html]
<app-counter [buttonText]="btnMsg"
             (counterClickedTenTimes)="addClicks($event)" />

<ul>
  @for (click of clicksList; track click) {
    <li>Il pulsante e' stato premuto {{click}} volte.</li>
  }
</ul>
\end{lstlisting}
\risultato
\begin{immagine}
\begin{tikzpicture}
  \node[wtnode, font=\FiraCodeNL\scriptsize, minimum width=4.4cm] (p) at (0,0)
    {AppComponent};
  \node[wtnodeacc, font=\FiraCodeNL\scriptsize, minimum width=4.4cm] (f) at (0,-2.5)
    {CounterComponent};
  \draw[wtarrow] ([xshift=-1.3cm]p.south) -- node[mnota, left, align=right]
    {\code{[buttonText]="btnMsg"}\\[2pt]\textbf{dati}: scendono} ([xshift=-1.3cm]f.north);
  \draw[wtarrowacc] ([xshift=1.3cm]f.north) -- node[mnota, right]
    {\code{emit(10)}\\\code{addClicks(10)}\\[2pt]\textbf{eventi}: salgono}
    ([xshift=1.3cm]p.south);
\end{tikzpicture}
\end{immagine}
\end{affianco}

\begin{affianco}[0.44][c]
\begin{lstlisting}[style=ts, name=app.component.ts]
export class AppComponent {
  name = 'Web Technologies';
  btnMsg = "CLICK HERE";
  clicksList: number[] = [];

  addClicks(nClicks: number) {
    this.clicksList.push(nClicks);
  }
}
\end{lstlisting}
\risultato
\begin{browser}[dopo 20 clic]
  \wtbottone{{\normalsize\Lato CLICK HERE}}
  \begin{itemize}[leftmargin=14pt, itemsep=0pt, topsep=4pt]
    \item Il pulsante e' stato premuto 10 volte.
    \item Il pulsante e' stato premuto 20 volte.
  \end{itemize}
\end{browser}
\end{affianco}

Si noti la simmetria della sintassi: \textbf{parentesi quadre} per i dati che
scendono dal padre al figlio, \textbf{parentesi tonde} per gli eventi che
salgono dal figlio al padre.

\section{Il Router}

Le Single Page Application devono gestire il routing lato client. Angular
include un \term{Router} che abilita la navigazione fra viste diverse,
interpretando la URL come un'istruzione per cambiare vista. Le applicazioni
create con \code{ng new} hanno già il routing configurato e pronto.

\subsection{Il punto di ingresso}

\begin{lstlisting}[style=ts, name=src/main.ts]
import { bootstrapApplication } from '@angular/platform-browser';
import { appConfig } from './app/app.config';
import { AppComponent } from './app/app.component';

bootstrapApplication(AppComponent, appConfig)
  .catch((err) => console.error(err));
\end{lstlisting}

\begin{affianco}[0.46][c]
\begin{lstlisting}[style=ts, name=src/app/app.config.ts]
import { ApplicationConfig } from '@angular/core';
import { provideRouter } from '@angular/router';
import { routes } from './app.routes';

export const appConfig: ApplicationConfig = {
  providers: [provideRouter(routes)]
};
\end{lstlisting}
\risultato
\begin{immagine}
\begin{tikzpicture}[
    f/.style={wtnodeplain, font=\FiraCodeNL\scriptsize},
    a/.style={-{Stealth[length=4pt,width=3.5pt]}, line width=0.75pt,
      draw=neutral!45!black}]
  \node[f] (m) at (0,0) {main.ts};
  \node[f] (c) at (-1.7,-1.3) {app.component.ts};
  \node[f] (k) at (1.7,-1.3) {app.config.ts};
  \node[f] (r) at (1.7,-2.6) {app.routes.ts};
  \draw[a] (m) -- node[mnota, left=2pt]{radice} (c);
  \draw[a] (m) -- node[mnota, right=2pt]{configurazione} (k);
  \draw[a] (k) -- node[mnota, right]{\code{provideRouter}} (r);
\end{tikzpicture}
\end{immagine}
\end{affianco}

\subsection{Le rotte}

\begin{affianco}[0.6][c]
\codicepiccolo
\begin{lstlisting}[style=ts, name=src/app/app.routes.ts]
import { Routes } from '@angular/router';
import { CounterPageComponent } from './counter-page/counter-page.component';
import { HomePageComponent } from './home-page/home-page.component';

export const routes: Routes = [
  {
    path: "",                     // la URL della rotta
    title: "Home Page",           // il titolo della scheda
    component: HomePageComponent  // il componente da disegnare
  },
  {
    path: "counter",
    title: "Counter",
    component: CounterPageComponent
  }
];
\end{lstlisting}
\risultato
\begin{immagine}
\begin{tikzpicture}[
    u/.style={mcod, draw=neutral!55!black, fill=white, line width=0.7pt,
      minimum width=1.5cm, minimum height=0.5cm},
    a/.style={-{Stealth[length=4pt,width=3.5pt]}, line width=0.75pt,
      draw=neutral!45!black}]
  \node[mnota, font=\Lato\scriptsize\bfseries] at (0,0.6) {URL};
  \node[mnota, font=\Lato\scriptsize\bfseries] at (3.3,0.6) {componente};
  \node[u] (u1) at (0,0) {/};
  \node[mfun] (c1) at (3.3,0) {HomePageComponent};
  \node[u] (u2) at (0,-1.3) {/counter};
  \node[mfun] (c2) at (3.3,-1.3) {CounterPageComponent};
  \draw[a] (u1) -- (c1); \draw[a] (u2) -- (c2);
  \node[mnota, anchor=north] at ([yshift=-1pt]c1.south) {scheda: «Home Page»};
  \node[mnota, anchor=north] at ([yshift=-1pt]c2.south) {scheda: «Counter»};
\end{tikzpicture}
\end{immagine}
\end{affianco}

\subsection{router-outlet e routerLink}

Il componente \code{RouterOutlet} analizza la URL corrente e disegna il
componente associato al percorso corrispondente.

\begin{affianco}[0.48][c]
\codicepiccolo
\begin{lstlisting}[style=ts, name=app.component.ts]
import { Component } from '@angular/core';
import { CommonModule } from '@angular/common';
import { RouterLink, RouterOutlet } from '@angular/router';

@Component({
  selector: 'app-root',
  standalone: true,
  imports: [CommonModule, RouterOutlet, RouterLink],
  template: `
    <nav>
      <a routerLink="/">Home</a> |
      <a routerLink="/counter">Counter</a>
    </nav>
    <router-outlet></router-outlet>
  `,
})
export class AppComponent {}
\end{lstlisting}
\risultato
\begin{browser}[localhost:4200/counter]
  \wtlink{Home} | \wtlink{Counter}\par\vspace{6pt}
  \tikz\node[draw=accent!60!black, dashed, rounded corners=2pt, inner sep=6pt,
    minimum width=0.9\linewidth, minimum height=1.4cm, font=\Lato\footnotesize,
    text=neutral!40!black,
    label={[font=\FiraCodeNL\tiny, text=accent!70!black, inner sep=1pt,
            anchor=south east]north east:<router-outlet>}]
    {\textit{qui viene disegnato} \code{CounterPageComponent}};
\end{browser}
\wtnota{La barra resta, il contenuto dell'outlet cambia con la URL.}
\end{affianco}

\error{Non usare \attr{href} nei collegamenti interni}{
  Usando collegamenti tradizionali con l'attributo \attr{href}, il browser
  effettua una nuova richiesta HTTP e \textbf{ricarica l'intera applicazione}.
  Il Router analizzerebbe poi la URL e caricherebbe il componente giusto, quindi
  il risultato finale sarebbe corretto --- ma vanificando completamente lo scopo
  di una Single Page App.

  Per navigare correttamente fra le viste si usa la direttiva
  \code{routerLink}. Si noti che, se un utente inserisce a mano la URL
  \code{localhost:4200/counter}, il caricamento completo avviene comunque e il
  Router disegna correttamente il componente: è il comportamento desiderato.}

\subsection{Evidenziare il collegamento attivo}

La direttiva \code{RouterLinkActive} permette di indicare un elenco di classi
CSS da applicare ai \code{routerLink} che corrispondono alla rotta corrente.

\begin{affianco}[0.5][c]
\codicepiccolo
\begin{lstlisting}[style=html, name=app.component.html]
<nav>
  <a routerLink="/" routerLinkActive="active"
     [routerLinkActiveOptions]="{exact: true}">Home</a> |
  <a routerLink="/counter" routerLinkActive="active">Counter</a>
</nav>
<router-outlet></router-outlet>
\end{lstlisting}

\begin{lstlisting}[style=css]
nav .active {
  font-weight: bold;
}
\end{lstlisting}
\risultato
\begin{browser}[localhost:4200]
  \textbf{\wtlink{Home}} | \wtlink{Counter}
\end{browser}
\begin{browser}[localhost:4200/counter]
  \wtlink{Home} | \textbf{\wtlink{Counter}}
\end{browser}
\end{affianco}

L'opzione \code{exact: true} serve per la rotta radice: senza di essa,
\code{/} risulterebbe attiva anche quando ci si trova su \code{/counter}, dato
che ogni percorso comincia con una barra.

\section{Riferimenti}

\begin{itemize}
  \item \textbf{Angular Docs}. \\
        \urlx{https://angular.dev/overview} \\
        Parti rilevanti: \textit{Introduction: Essentials}; \textit{In-depth
        Guides: Components, Template Syntax, Routing}; \textit{Developer Tools:
        Angular CLI}.
\end{itemize}
